\chapter*{Introduction Générale}
\addcontentsline{toc}{chapter}{Introduction Générale}
\markboth{Introduction Générale}{Introduction Générale}

Le secteur du \textbf{BTP} (\textit{Bâtiment et Travaux Publics}) exige une coordination
rigoureuse entre acteurs, documents et échéances. Chaque chantier mobilise des équipes
pluridisciplinaires, génère des devis, contrats, factures et bons de livraison, et nécessite
un suivi continu de l'avancement. Pourtant, de nombreuses entreprises s'appuient encore sur
des outils dispersés --- tableurs, documents papier, échanges informels --- qui fragmentent
l'information, ralentissent la prise de décision et augmentent le risque d'erreurs. La
digitalisation, notamment via un \textbf{ERP} (\textit{Enterprise Resource Planning}) adapté
au BTP, permet de centraliser la gestion des projets, des clients et des documents commerciaux
au sein d'une plateforme unique.

\medskip

Ce projet de fin d'études, réalisé au sein de \textbf{CIVCO BTP}, consiste à concevoir et
développer une \textbf{application web} de gestion et de planification des projets de
construction de l'entreprise. Cette solution vise à accompagner les équipes internes dans le
pilotage quotidien de leurs chantiers : suivi des projets et des tâches, gestion des clients,
production des documents commerciaux, suivi des contrats et consultation d'un tableau de bord
de synthèse. Elle répond ainsi au besoin de fiabiliser les processus métier et d'améliorer la
collaboration entre les différents services.

\medskip

La solution repose sur une architecture \textit{full-stack} : un \textit{backend} API
\textbf{Laravel} / \textbf{PHP} exposant des services \textbf{REST} sécurisés, couplé à une
interface \textbf{React} en \textbf{SPA} (\textit{Single Page Application}), avec un contrôle
d'accès par rôles (\textbf{RBAC}) adapté aux profils de l'entreprise.

\medskip

Ce rapport est organisé en quatre chapitres :
\begin{itemize}
  \setlength{\itemsep}{0pt}
  \item \textbf{Chapitre 1} : présentation de CIVCO BTP, analyse de la problématique,
  objectifs du projet et méthodologie adoptée.
  \item \textbf{Chapitre 2} : état de l'art et justification des choix technologiques
  (\textit{full-stack}, Laravel, React, ERP BTP).
  \item \textbf{Chapitre 3} : conception du système (besoins fonctionnels, acteurs, architecture,
  modélisation \textbf{UML}).
  \item \textbf{Chapitre 4} : réalisation des modules principaux et validation de
  l'application par des tests.
\end{itemize}

\newpage
